% Chapter 3: Cipher Maps
% DEFINES: ch:cipher-maps; sec:machines, sec:tuple, sec:three-transformations,
%   sec:totality-support, sec:cipher-type, sec:ch3-notes;
%   def:cipher-map, ex:watchlist, rem:values-are-maps
% REFERENCES (backward, Part I): sec:oblivious-service, sec:untrusted-sees
% REFERENCES (forward, footnote part labels): part:algebra, part:confidentiality,
%   part:constructions

\chapter{Cipher Maps}
\label{ch:cipher-maps}

Part~I described, in words, a machine computing on encodings it cannot
read. This chapter makes the object precise. By its end the cipher map is a
formal tuple, totality is a stated property rather than a promise, the three
transformations that produce a cipher map from an ordinary function are laid
out, and the support against which the next chapter measures leakage is
named. The definitions here are the book's definitions of record; the later
parts add structure to them but do not change them.

\section{The Trusted and Untrusted Machines}
\label{sec:machines}

Fix a \emph{latent function} $f : X \to Y$ whose domain, codomain, and rule
are known only to a trusted party. The arrangement of
\Cref{sec:oblivious-service} has two machines. The \emph{trusted machine},
$T$, holds an encoder $\enc$, a decoder $\dec$, and a secret $s$; it is the
only party that can move between values and their encodings. The
\emph{untrusted machine}, $U$, holds a single function $\fhat$ on bit
strings and evaluates it on whatever it is given. To evaluate $f$ at $x$,
the trusted machine encodes $x$ as a token $c = \enc(x)$, the untrusted
machine returns $\fhat(c)$, and the trusted machine decodes the result as
$\dec(\fhat(c))$ (\Cref{fig:two-machines}).

\begin{figure}[t]
\centering
\begin{tikzpicture}[
  box/.style={draw, rounded corners, align=center, minimum height=1.3cm,
              minimum width=3.6cm, font=\small},
  >={Stealth[]}]
\node[box] (T) {Trusted machine $T$\\[2pt]{\footnotesize holds $\enc,\dec,s$}};
\node[box] (U) at (7.2,0) {Untrusted machine $U$\\[2pt]{\footnotesize holds $\fhat$}};
\draw[->] ([yshift=3mm]T.east) -- node[above,font=\footnotesize]{$c=\enc(x)$}
          ([yshift=3mm]U.west);
\draw[->] ([yshift=-3mm]U.west) -- node[below,font=\footnotesize]{$\fhat(c)$}
          ([yshift=-3mm]T.east);
\end{tikzpicture}
\caption{The two machines. The trusted machine encodes an input and decodes
the returned token; the untrusted machine evaluates the total function
$\fhat$ on the token it is given. The secret $s$ never leaves $T$, and the
result the trusted machine reads is $\dec(\fhat(c))$.}
\label{fig:two-machines}
\end{figure}

The division of labor is the whole point, so it is worth stating what each
machine can and cannot do. The trusted machine is small and holds the
secret; it does the encoding and decoding, which are cheap, and it never
ships the secret anywhere. The untrusted machine is large and holds the
table; it does the bulk computation, the lookup over a structure that may
encode an entire corpus, and it does so on tokens whose meaning it has no
way to recover. Nothing in the protocol asks the untrusted machine to be
honest about anything except returning $\fhat$ of what it is given, and
even a dishonest server that logs every token and studies them at leisure
is the adversary the rest of the book quantifies against.

\begin{example}[A watchlist]
\label{ex:watchlist}
A small running example will carry the abstraction through this part. Let
the universe of names be $X = \{\mathrm{alice}, \mathrm{bob},
\mathrm{carol}, \mathrm{dave}\}$ and let the latent function be membership
in a watchlist $S = \{\mathrm{alice}, \mathrm{carol}\}$, so $f(x) = 1$ if
$x \in S$ and $0$ otherwise. The trusted machine wants an untrusted server
to answer ``is this name on the list?'' without learning the name, the
list, or even how often a given name is queried. Concretely, to check
alice the trusted machine computes a token $c = \enc(\mathrm{alice})$, an
$n$-bit string that looks like a random draw; the server returns $\fhat(c)$,
another $n$-bit string; and the trusted machine decodes it to $1$. The
server has performed a membership test over an encrypted list and learned
neither the query nor the answer. We return to this map as each property is
added.
\end{example}

\section{The Cipher Map Tuple}
\label{sec:tuple}

\begin{definition}[Cipher map]
\label{def:cipher-map}
A \emph{cipher map} for a latent function $f : X \to Y$ is a tuple
$(\fhat, \enc, \dec, s)$ where
\begin{itemize}
\item $\fhat : \B^n \to \B^n$ is a \emph{total} function on $n$-bit strings,
\item $\enc : X \times \{0, \ldots, K(x) - 1\} \to \B^n$ encodes an input
  $x$ as one of $K(x) \geq 1$ tokens,
\item $\dec : \B^n \to Y \cup \{\bot\}$ decodes a token to an output value
  or to the symbol $\bot$ (``no value''),
\item $s$ is a secret, the trapdoor, from which $\fhat$, $\enc$, and $\dec$
  are derived.
\end{itemize}
The untrusted machine holds $\fhat$; the trusted machine holds $\enc$,
$\dec$, and $s$.
\end{definition}

Three features of the definition deserve a closer look, because each is
load-bearing later. The first is the type of $\fhat$: it is a function from
$\B^n$ to $\B^n$, bit strings to bit strings, with no mention of $X$ or
$Y$. The untrusted machine never sees the latent types; it sees an endo
function on tokens. The second is the second argument of $\enc$: a value
$x$ may be encoded as any of $K(x)$ distinct tokens, and which one is used
on a given call is the trusted machine's choice. This freedom is dead
weight in an ordinary encoding, but here it is the raw material for hiding
frequency, since a value that is queried often can be spread across many
tokens so that no single token is queried often. The third is the $\bot$ in
the range of $\dec$: not every token decodes to a value, and the fraction
that do is a parameter the construction controls.

The encodings are realized through a cryptographic hash $h : \B^{*} \to
\B^n$, which we model as a random oracle when uniformity of the tokens is
needed. The notation recurs throughout the book: $\B = \{0,1\}$, so $\B^n$
is the token space; $\fhat$ and $\ghat$ are cipher-map approximations of
latent functions $f$ and $g$; $\enc$ and $\dec$ are the trusted machine's
private operations; and the secret $s$ keys all three. For the watchlist,
$X$ is the four names, $Y = \{0,1\}$, and $\enc(\mathrm{alice}, k)$ is the
$k$-th of $K(\mathrm{alice})$ tokens for alice.

\begin{remark}[Values are constant maps]
\label{rem:values-are-maps}
A cipher value $\enc(v, k)$ is the same kind of object as a cipher map: it
is a cipher map for the constant function $f_v(x) = v$. Under this view the
untrusted machine deals with one kind of thing, total functions on bit
strings, some of which happen to be constant. This collapse of values and
functions is not a notational convenience only; it is what lets a cipher
map's output be fed as another's input, and it is the seed of the algebra
of \Cref{part:algebra}.\footnote{The cipher type $\cipher{X}$, the type of
cipher values over $X$, and the typed composition of cipher maps are
developed in \Cref{part:algebra}.}
\end{remark}

\section{Three Transformations}
\label{sec:three-transformations}

A cipher map is not built in one step, and seeing it as a composition of
three independent transformations explains why it has the properties it
has. Start with an ordinary function $f : X \to Y$ and apply, in turn:

\begin{description}
\item[Undefined injection.] Extend the domain $X$ to $X \cup \{\bot_1,
  \ldots, \bot_N\}$ by adding $N$ undefined elements, lifting $f$ to a
  partial function that is undefined on the $\bot_i$. The real inputs are
  now a fraction $\lvert X \rvert / (\lvert X \rvert + N)$ of an enlarged
  domain. This dilution is what gives filler queries somewhere to live: a
  decoy is just a draw from the undefined part. The dilution ratio is the
  parameter $\varepsilon$, the probability that an arbitrary token is
  treated as meaningful.
\item[Noise closure.] Lift the partial function to a \emph{total} one by
  sending every undefined input to a hash-determined output rather than to
  an error. After this step there are no undefined inputs to observe: every
  token produces an answer. This is the transformation that yields
  totality, and with it the indistinguishability of real and filler
  queries.
\item[Multiple representations.] Replace each value $x$ with $K(x) \geq 1$
  distinct tokens, keyed by the secret. Representational equality implies
  value equality, but not conversely: two tokens for the same value are
  unlinkable without the secret. Choosing $K(x)$ in proportion to how often
  $x$ is queried flattens the token distribution, so frequent values do not
  betray themselves. This is the transformation behind representation
  uniformity.
\end{description}

\Cref{tab:transformations} collects the three.

\begin{table}[t]
\centering
\caption{The three transformations and what each produces.}
\label{tab:transformations}
\begin{tabular}{@{}lll@{}}
\toprule
Transformation & Operation & Produces \\
\midrule
Undefined injection & extend $X$ with $N$ undefined elements & $\varepsilon$ (noise-decode prob.) \\
Noise closure & undefined inputs $\mapsto$ hash outputs & totality \\
Multiple representations & each value $\mapsto K(x)$ tokens & representation uniformity \\
\bottomrule
\end{tabular}
\end{table}

The four properties of the next chapter are then not a list to be
memorized but consequences of this construction: totality comes from noise
closure, representation uniformity from the multiple-representation step,
the noise-decode probability $\varepsilon$ from undefined injection, and the
remaining two, correctness and composability, from how the secret is chosen
and how cipher maps are chained. We present the three as conceptual layers,
not as algebraic monads; whether they satisfy formal monad laws in the
approximate setting is a question we neither need nor try to settle. The
decomposition earns its place by explaining the origin of each property,
and a reader who prefers to take \Cref{def:cipher-map} as given may do so
and skip ahead.

There are two ways to perform the construction, and the difference matters
in Part~V. In a \emph{batch} construction the trusted machine searches for
a secret $s$ that makes $\fhat$ correct on all of $X$ at once, paying
computation up front in exchange for tunable error and near-optimal space;
this requires knowing $X$ and $f$ at build time, as the HashSet and entropy
map do. In an \emph{online} construction $\fhat$ is read directly off a
hash structure with no search, the secret is just a hash key, and the
operations are limited to those the structure supports algebraically, as in
the trapdoor Boolean algebra. Both produce objects satisfying
\Cref{def:cipher-map}; they differ in cost, in which operations they admit,
and in how their parameters trade off.\footnote{The concrete batch and
online constructions, and their space and accuracy, are the subject of
\Cref{part:constructions}.}

\section{Totality and the Populated Support}
\label{sec:totality-support}

Return to the word \emph{total} in \Cref{def:cipher-map}. The function
$\fhat$ is defined on all $2^n$ bit strings, not only on the tokens that
encode real inputs. Hand it a string that encodes nothing and it still
returns an $n$-bit answer; there is no exception, no ``not found,'' no
out-of-band signal that the input was meaningless. This is the structural
fact from which the framework's first privacy guarantee follows: with no
special case to observe, the untrusted machine has no way to separate
inputs it should take seriously from inputs it should not.

Two quantities organize what the untrusted machine sees. The first is the
\emph{noise-decode probability} $\varepsilon$, already met as the dilution
ratio of undefined injection: it is the probability that a uniformly random
token decodes to a value rather than to $\bot$. A construction sets
$\varepsilon$ by how much of the $2^n$ token space it treats as meaningful,
and it pays for a small $\varepsilon$ in bits, roughly $-\log_2 \varepsilon$
of them per element, because rejecting more strings as $\bot$ takes a
sparser meaningful set. The second is the \emph{populated support}
$\mathrm{im}(\enc) \subseteq \B^n$, the set of tokens that actually encode
real inputs. Its size $\lvert \mathrm{im}(\enc) \rvert$ is small relative to
$2^n$, and the observed token distribution lives on it, not on all of
$\B^n$. This distinction looks pedantic and is not: the next chapter
measures confidentiality as closeness to the uniform distribution \emph{on
this support}, and measuring against the uniform distribution on all of
$\B^n$ instead would make every sparse construction look maximally leaky and
the measure say nothing. Getting the support right is the difference between
a meaningful confidentiality measure and a vacuous one.

\begin{example}[The watchlist, in numbers]
\label{ex:watchlist-numbers}
Take $n = 12$, so a token is one of $2^{12} = 4096$ bit strings. Suppose
alice is queried often enough to warrant $K(\mathrm{alice}) = 4$
representations and the other three names two each, so the populated
support $\mathrm{im}(\enc)$ holds $4 + 2 + 2 + 2 = 10$ tokens, a $0.24\%$
slice of the token space. To check alice the trusted machine picks one of
her four tokens, say $\enc(\mathrm{alice}, 1) = \mathtt{0x2A7}$ (the string
$0010\,1010\,0111$); the server returns $\fhat(\mathtt{0x2A7})$, say
$\mathtt{0x6B4}$; the trusted machine decodes $\mathtt{0x6B4}$ to $1$, on
the list. Two small sets live inside the $4096$, at the two ends of the
pipeline $\dec(\fhat(\enc(x,k)))$, and they play different roles. The
\emph{input} support is the ten tokens just counted, the ones that
\emph{encode} the names; it is what the next chapter's $\delta$ measures
uniformity against, and it fixes $H^* = \log_2 10$. The \emph{output}
acceptance set is separate: hand $\dec$ a uniformly random $12$-bit string
and it returns a $0$ or $1$ rather than $\bot$ only with probability
$\varepsilon$, the fraction of strings that \emph{decode} to an answer; if
the construction reserves four of the $4096$ codewords as valid answers,
then $\varepsilon \approx 4/4096 \approx 10^{-3}$ and an out-of-domain probe
is rejected about $999$ times in $1000$. The ten input tokens govern
$\delta$; the four output codewords govern $\varepsilon$; they are different
sets, and neither is the $4096$ ambient strings.
\end{example}

Totality also makes filler queries free. Because every string produces
output, the trusted machine can interleave random tokens among the real
ones, and the untrusted machine cannot tell which are which. In the
watchlist, the trusted machine can issue a stream of lookups of which some
fraction are decoys drawn from the undefined part of the domain, and the
server, seeing only that each token yields an answer, cannot subtract the
decoys back out. Totality is the property that the four cannots of
\Cref{sec:untrusted-sees} called ``cannot distinguish a real query from
filler,'' now stated as a structural feature of $\fhat$ rather than a hope,
and the next chapter makes it the first of the four.

\section{The Cipher Type}
\label{sec:cipher-type}

It will be useful, and it costs little here, to name the type-level view
that \Cref{part:algebra} develops in full. Write $\cipher{X}$ for the
\emph{cipher type} over $X$: the type whose values are tokens encoding
elements of $X$. A cipher map for $f : X \to Y$ is then a morphism that
sends $\cipher{X}$ to $\cipher{Y}$, an arrow at the level of cipher types
rather than of plain values, and \Cref{rem:values-are-maps} is the
statement that the objects and the morphisms of this picture are made of
the same material: a cipher value in $\cipher{X}$ is a cipher map from a
trivial domain.

This is the vantage from which the algebra of cipher maps becomes visible.
If cipher maps are arrows between cipher types, one can ask which type
constructors, product, sum, function, survive the passage through the
trapdoor and which do not, and the answer turns out to be sharply
different for different constructors: product structure passes cleanly,
sum structure runs into an impossibility. That is the subject of
\Cref{part:algebra}; here it is enough to have the word $\cipher{X}$ and
the picture of cipher maps as arrows, so that the four properties of the
next chapter can be read as properties of those arrows.

\section{Notes and Provenance}
\label{sec:ch3-notes}

\Cref{def:cipher-map} is the cipher map of the formal specification and of
the cipher-maps paper; the tuple, the totality requirement, and the
trusted/untrusted split are taken from there without change. The reading of
the construction as three orthogonal transformations of a type
(\Cref{sec:three-transformations}) comes from the algebraic-cipher-types
design work; it is the conceptual reason each property exists and is not
needed to use the definitions. The batch and online strategies, and the
HashSet, entropy map, and trapdoor Boolean algebra named as their examples,
are developed in \Cref{part:constructions}. The view of cipher values as
constant cipher maps (\Cref{rem:values-are-maps}) is the hinge that
\Cref{part:algebra} turns into a typed algebra, and the choice to measure
confidentiality against the populated support rather than the ambient token
space (\Cref{sec:totality-support}) is the point on which the confidentiality
theory of \Cref{part:confidentiality} is built.
